\chapter{Bloco 0 -- Por que programação inteira? (aquecimento)}
\label{cap:bloco0}

Este bloco revisita os exemplos do Capítulo~\ref{cap:arredondar}, agora no código. Há um único exercício para
completar (o modelo de Wolsey); o resto são células para rodar e interpretar.

\section{0.1 -- Exercício: o modelo de Wolsey em PuLP}

\begin{discussao}
\begin{itemize}
  \item \emph{Qual a ideia por trás de ``resolver o PL e arredondar''?} Aproveitar a eficiência do simplex,
  supondo que o ótimo inteiro esteja ``perto'' do ótimo contínuo.
  \item \emph{O ótimo do PL será inteiro?} Não: é $(1{,}95;\,4{,}92)$. \emph{O arredondamento será viável?}
  $(2,5)$ não é. \emph{Será ótimo?} O único vizinho viável, $(2,4)$, vale $4{,}56$, e o ótimo inteiro é $5$ em
  $(5,0)$. Veja a Seção~\ref{sec:wolsey}.
\end{itemize}
\end{discussao}

\paragraph{O que o exercício pede.} A função já cria o problema e as duas variáveis; o tipo das variáveis depende do
parâmetro \py{inteiro}. Faltam três linhas: o objetivo e as duas restrições.

\paragraph{Passo a passo.}
\begin{enumerate}
  \item Objetivo $\max\; x_1 + 0{,}64x_2$: como o problema foi criado com \py{LpMaximize}, basta somar a expressão.
  \item Restrição $50x_1+31x_2\le250$: expressão com \py{<=}.
  \item Restrição $3x_1-2x_2\ge-4$: expressão com \py{>=}. Não é preciso ``passar para $\le$'': o PuLP aceita
  \py{>=} diretamente.
\end{enumerate}

\begin{codigo}
def modelo_wolsey(inteiro=True):
    cat = pulp.LpInteger if inteiro else pulp.LpContinuous
    prob = pulp.LpProblem("wolsey", pulp.LpMaximize)
    x1 = pulp.LpVariable("x1", lowBound=0, cat=cat)
    x2 = pulp.LpVariable("x2", lowBound=0, cat=cat)
    prob += 1.00*x1 + 0.64*x2, "obj"        # função objetivo
    prob += 50*x1 + 31*x2 <= 250, "c1"      # restrição 1
    prob += 3*x1 - 2*x2 >= -4, "c2"         # restrição 2
    return prob, x1, x2
\end{codigo}

\paragraph{Verificação.} A célula roda o modelo duas vezes, com \py{inteiro=False} (PL) e com \py{inteiro=True}
(PI), usando a função \py{rodar}. Os valores esperados são $5{,}0984$ e $5$. A célula seguinte desenha a
Figura~\ref{fig:wolsey} e lista os quatro vizinhos inteiros do ótimo da PL, como na tabela da
Seção~\ref{sec:wolsey}.

\begin{atencao}
\begin{itemize}
  \item Escrever o objetivo depois das restrições funciona, mas ele precisa ser a única expressão
  \emph{sem} comparação somada ao problema.
  \item Se aparecer ``Falta definir a função objetivo'', a linha do objetivo não foi executada. Lembre-se de rodar a
  célula da função \emph{antes} da célula de verificação.
\end{itemize}
\end{atencao}

\section{0.1 (continuação) -- O caso binário}

A célula da mochila de 3 itens não tem \py{TODO}: ela enumera todas as $2^3=8$ soluções com
\py{itertools.product([0, 1], repeat=3)}, que gera todas as tuplas $(x_1,x_2,x_3)$ com valores 0 ou 1. Para cada uma,
calcula peso e valor e marca se é viável. A saída mostra o que vimos na Seção~\ref{sec:mochila3}: a gulosa (razões
$4$, $3{,}5$, $3$) pega o item 1 e fica em 12; o ótimo é 13, com os itens 2 e 3.

\begin{resumo}[O que levar desta célula]
Enumerar funciona para $n=3$ (8 casos), mas $2^n$ explode muito rápido. Gulosas são rápidas, mas sem garantia. É
por isso que precisamos de modelos e solvers.
\end{resumo}

\section{0.2 -- Escala de enfermeiras}

Também não há \py{TODO}, mas vale a pena ler o código com cuidado, porque ele usa padrões que você vai repetir.

\begin{codigo}
d = [17, 13, 15, 19, 14, 16, 11]          # demanda de dom..sáb

def modelo_enfermeiras(d, inteiro=True):
    prob = pulp.LpProblem("enfermeiras", pulp.LpMinimize)
    cat = pulp.LpInteger if inteiro else pulp.LpContinuous
    x = [pulp.LpVariable(f"x_{k}", lowBound=0, cat=cat) for k in range(7)]
    prob += pulp.lpSum(x), "total"
    for j in range(7):
        prob += pulp.lpSum(x[k] for k in range(7) if (j - k) % 7 <= 4) >= d[j], f"dia_{j}"
    return prob, x
\end{codigo}

\begin{itemize}
  \item \py{x[k]} = enfermeiras que \textbf{começam} no dia $k$ (veja a Seção~\ref{sec:enfermeiras} para entender por
  que essa é a variável certa).
  \item \py{pulp.lpSum(x)} soma todas as variáveis da lista: o total contratado.
  \item O laço cria uma restrição por dia $j$. O gerador soma apenas os $x_k$ cujo turno de 5 dias inclui o dia $j$,
  isto é, $(j-k)\bmod 7\in\{0,1,2,3,4\}$. O operador \texttt{\%} do Python devolve sempre um resto entre 0 e 6, mesmo
  quando $j-k$ é negativo, o que faz a ``volta'' de sábado para domingo funcionar sozinha.
\end{itemize}

\paragraph{Saída e leitura.} PL $=22{,}33$, PI $=23$, PL arredondada para cima $=25$. A mensagem central:
\textbf{arredondar para cima garante viabilidade, mas não otimalidade}; e o limite da PL ($22{,}33$) prova que 23 é o
melhor possível, já que o total é inteiro e precisa ser $\ge 22{,}33$.

\section{O mapa do resto da aula}

A partir do Bloco 1, a inteireza deixa de vir da natureza do problema e passa a ser uma \textbf{ferramenta de
modelagem}. Cada bloco ensina uma ``receita'' de linearização, e todas estão reunidas no
Capítulo~\ref{cap:resumo}:

\begin{center}
\begin{tabular}{cl}
\toprule
Bloco & Técnica\\
\midrule
1 & Implicações \emph{se--então} (custo fixo, produção condicionada)\\
2 & Restrições disjuntivas (\emph{big-M}, escalonamento)\\
3 & Relações lógicas (tabela-verdade, carteira de investimentos)\\
4 & Representação de valores discretos\\
5 & ``Ou exclusivo'' e indicadores (teoria de variáveis lógicas)\\
6 & Problema da mochila binária\\
7 & \emph{Bin packing}\\
\bottomrule
\end{tabular}
\end{center}
